\documentclass[12pt]{article}
\usepackage{acronymtooltips}
% Your acronym definitions — edit this file only.
% Use \xx{key} and \xxx{key} in your document body.

\newacronym{bmof}{BMOF}{Block Matching Optical Flow}
\newacronym{bm}{BM}{Block Matching}
\newacronym[description={Multipurpose Block RAM memory module in FPGA}]{bram}{BRAM}{Block RAM}
\newacronym{dnn}{DNN}{Deep Neural Network}
\newacronym{dram}{DRAM}{Dynamic RAM}
\newacronym[description={Direction Selective model of motion detection in biological vision, usually either Hassenstein-Reichhardt or Barlow-Levick type}]{ds}{DS}{Direction Selective}
\newacronym{fpga}{FPGA}{Field Programmable Gate Array}
\newacronym{gt}{GT}{Ground Truth}
\newacronym{of}{OF}{Optical Flow}
\newacronym[description={Register Transfer Logic intermediate form, consisting of combinational and synchronous register logic cells}]{rtl}{RTL}{Register Transfer Logic}
\newacronym{timsl}{TS}{time slice}


\begin{document}

\setlength{\parskip}{\baselineskip}%

This sample project demonstrates how you can get tooltip popups for your acronyms, along with highlighting in bold and color the original acronym definition.

Upload \texttt{acronymtooltips.sty}, \texttt{acronymtooltips-pdfcomment.tex}, and \texttt{myacronyms.tex} to Overleaf, then use
\texttt{\string\usepackage\{acronymtooltips\}} and \texttt{\string\input\{myacronyms\}} in your preamble.

Use \texttt{\string\xx\{ac\}} for singular form and \texttt{\string\xxx\{ac\}} for plural form.

See \textit{myacronyms.tex} for your acronym definitions and \textit{acronymtooltips.sty} for colors and tooltip logic.

Here is some sample text:
Testing \xx{of} algorithms is hard because \xx{of} requires good \xx{gt} data.
Also, you need lots of \xx{bram} because \xx{of} takes a lot of memory, particularly for \xx{bm}.
\xx{bm} produces flow vectors and can be combined with corner detection to disambiguate flow along edges.
It is much more useful that simplified biologically-inspired \xx{ds} flow.
But \xx{bram} is costly.
Using many \xxx{bram}
will cost you a lot of money for big \xxx{fpga}.
But it enabled \xx{bm} \xx{of} on \xxx{fpga}, more accurate than \xx{ds} models of the past and results much closer to \xx{gt} flow.

\end{document}
